\subsection{Asymptotes and global behavior}

Limits describe not only local holes and jumps but also persistent geometric behavior.

\begin{definitionbox}[Vertical and horizontal asymptotes]
The line \(x=a\) is a vertical asymptote when at least one one-sided limit of \(f(x)\) as \(x\to a\) is infinite. The line \(y=L\) is a horizontal asymptote when
\[
\lim_{x\to\infty}f(x)=L
\quad\text{or}\quad
\lim_{x\to-\infty}f(x)=L.
\]
\end{definitionbox}

\begin{examplebox}[Two asymptotes from one simple function]
For
\[
f(x)=\frac{x+1}{x-1},
\]
the denominator vanishes at \(x=1\), while the numerator tends to \(2\), so \(x=1\) is a vertical asymptote. Dividing numerator and denominator by \(x\) gives
\[
\frac{1+1/x}{1-1/x}\to1,
\]
so \(y=1\) is a horizontal asymptote.
\end{examplebox}

\begin{interpretationbox}[An asymptote is a behavior statement]
An asymptote does not mean that a graph can never cross the line. It means that a particular limiting relationship holds near a boundary point or far from the origin.
\end{interpretationbox}
